\documentclass[10pt,a4paper]{article}
\usepackage{amsmath,amssymb,amsthm,graphicx}
\newif\iflipics\lipicsfalse
\usepackage{amsmath}
\usepackage{amsfonts}
\usepackage{amssymb}
\usepackage{amsthm}
\usepackage{xspace}
\usepackage[dvipsnames]{xcolor}
\usepackage{tabularx}
\usepackage{ragged2e}
\usepackage{nicefrac}
\usepackage{bbm}
\usepackage{caption}
\usepackage{graphicx}
\usepackage{enumitem}
\usepackage{soul}
\usepackage{hyperref}
\definecolor{darkgreen}{rgb}{0,0.4,0}
\usepackage{algorithm}
\usepackage[commentColor=darkgreen]{algpseudocodex}
\newcommand{\LineComment}{\LComment}
% \algrenewcomment[1]{\hfill \textit{\textcolor{darkgreen}{$\triangleright$ #1}}}
% \algnewcommand{\LineComment}[1]{\State \textit{\textcolor{darkgreen}{$\triangleright$ #1}}}
\algnewcommand{\IfThenElse}[3]{% \IfThenElse{<if>}{<then>}{<else>}
  \State \algorithmicif\ #1\ \algorithmicthen\ #2\ \algorithmicelse\ #3}
\algdef{SE}[SUBALG]{Indent}{EndIndent}{}{\algorithmicend\ }%
\makeatletter
\renewcommand{\ALG@name}{Listing}
\makeatother
\usepackage{cleveref}
\crefname{listing}{listing}{listings}
\crefname{equation}{eq.}{eqs.}

% must be after cleveref
\newtheorem{theorem}{Theorem}[section]
\newtheorem{lemma}[theorem]{Lemma}
\newtheorem{corollary}[theorem]{Corollary}
\newtheorem{fact}{Fact}[section]
\newtheorem{conj}[theorem]{Conjecture}
\newtheorem{question}[theorem]{Question}
\newtheorem{claim}[theorem]{Claim}
\newtheorem{obs}[theorem]{Observation}
\newtheorem{proposition}[theorem]{Proposition}
\theoremstyle{definition}
\newtheorem{definition}[theorem]{Definition}
\newtheorem{example}{Example}
\newtheorem{results}{Result Family}
\newtheorem{result}{Result}
\theoremstyle{remark}
\newtheorem*{remark}{Remark}
\makeatletter
\newcommand\squeezepar{\@startsection{paragraph}{4}%
{\z@}{1.5ex \@plus1ex \@minus.2ex}{-1em}{\normalfont\normalsize\bfseries}}
\newcommand{\mypara}[1]{\squeezepar{{#1}.}}
\renewcommand{\paragraph}[1]{\squeezepar{{#1}.}}
\makeatother
\graphicspath{{items/}{./}}
\newcommand{\EE}{\mathbb{E}}
\newcommand{\FF}{\mathbb{F}}
\newcommand{\NN}{\mathbb{N}}
\newcommand{\RR}{\mathbb{R}}
\newcommand{\ZZ}{\mathbb{Z}}
\newcommand{\indic}{\mathbbm{1}}
\newcommand{\cA}{\mathcal{A}}
\newcommand{\cB}{\mathcal{B}}
\newcommand{\cC}{\mathcal{C}}
\newcommand{\cD}{\mathcal{D}}
\newcommand{\cE}{\mathcal{E}}
\newcommand{\cF}{\mathcal{F}}
\newcommand{\cG}{\mathcal{G}}
\newcommand{\cH}{\mathcal{H}}
\newcommand{\cI}{\mathcal{I}}
\newcommand{\cJ}{\mathcal{J}}
\newcommand{\cM}{\mathcal{M}}
\newcommand{\cO}{\mathcal{O}}
\newcommand{\cP}{\mathcal{P}}
\newcommand{\cQ}{\mathcal{Q}}
\newcommand{\cR}{\mathcal{R}}
\newcommand{\cS}{\mathcal{S}}
\newcommand{\cU}{\mathcal{U}}
\newcommand{\cV}{\mathcal{V}}
\newcommand{\cW}{\mathcal{W}}
\newcommand{\cX}{\mathcal{X}}
\newcommand{\cY}{\mathcal{Y}}
\newcommand{\cZ}{\mathcal{Z}}

\newcommand{\tG}{\widetilde{G}}
\newcommand{\tS}{\widetilde{S}}
\newcommand{\tO}{\widetilde{O}}
\newcommand{\tP}{\widetilde{P}}
\newcommand{\tQ}{\widetilde{Q}}
\newcommand{\tf}{\widetilde{f}}
\newcommand{\tTheta}{\widetilde{\Theta}}
\newcommand{\tOmega}{\widetilde{\Omega}}
\newcommand{\tM}{\widetilde{M}}
\newcommand{\tmu}{\widetilde{\mu}}

\newcommand{\hw}{\hat{w}}
\newcommand{\hb}{\hat{b}}
\newcommand{\hP}{\hat{P}}
\newcommand{\hf}{\hat{f}}
\newcommand{\hQ}{\hat{Q}}

\newcommand{\ceil}[1]{{\left\lceil{#1}\right\rceil}}
\newcommand{\floor}[1]{{\left\lfloor{#1}\right\rfloor}}
\newcommand{\norm}[1]{{\left\Vert{#1}\right\Vert}}
\DeclareSymbolFont{newfont}{OML}{cmm}{m}{it}% Computer Modern math font
\DeclareMathSymbol{\emptystream}{3}{newfont}{15}% Symbol 15 "epsilon"
\let\epsilon\varepsilon
\renewcommand{\emptyset}{\varnothing}

\newcommand{\avoid}{\textsc{avoid}\xspace}
\newcommand{\mif}{\textsc{mif}\xspace}
\newcommand{\miflong}{\textsc{MissingItemFinding}\xspace}
\newcommand{\fco}{\textsc{fco}\xspace}
\newcommand{\fao}{\textsc{fao}\xspace}
\newcommand{\init}{\mathrm{init}\xspace}
\newcommand{\poly}{\mathrm{poly}\xspace}
\newcommand{\polylog}{\mathrm{polylog}\xspace}
\newcommand{\Algs}{\textsc{Algs}\xspace}

\newcommand{\repeatsc}{\textsc{repeat}\xspace}
\newcommand{\unsafesc}{\textsc{unsafe}\xspace}
\newcommand{\incompletesc}{\textsc{incomplete}\xspace}
\newcommand{\failsc}{\textsc{fail}\xspace}
\newcommand{\emptysc}{\textsc{empty}\xspace}
\newcommand{\timeoutsc}{\textsc{timeout}\xspace}
\newcommand{\bigsc}{\textsc{big}\xspace}
\newcommand{\abort}{\textsc{abort}\xspace}
\newcommand{\sort}{\textsc{sort}\xspace}
\newcommand{\abortsc}{\textsc{abort}\xspace}
\newcommand{\errorsc}{\textsc{error}\xspace}

\newcommand{\rl}{\mathbf{\ell}}             % stream length
\newcommand{\stream}{\tau}                  % generic input stream
\newcommand{\morestream}{\varphi}           % ``extension'' of \stream
\newcommand{\concat}{\mathbin{\circ}}       % concatenation of sequences
\newcommand{\trans}{T}                      % transition function of automaton
\newcommand{\SSD}[2]{\textsc{seqs}(#1,#2)}  % sorted sequences of distinct things
\newcommand{\Outs}{\textsc{Outs}\xspace}    % output seq given start state + adversary
\newcommand{\Mstk}{\Delta}                  % min max probability of mistake

\newcommand{\todo}[1]{%
  \marginpar{\color{red}!!!}[{\color{red}\bf  #1}]}
\newcommand{\gobble}[1]{}
\crefname{originalsection}{original section}{original sections}
\Crefname{originalsection}{Original section}{Original sections}
\makeatletter
\newlabel{subsec:rt-lb-overview}{{2.2}{source}{}{sourcepointer-1}{}}
\newlabel{subsec:rt-lb-overview@cref}{{[originalsection][0][]2.2}{[1][0][]source}{}{}{}}
\makeatother
\crefname{originalreference}{original reference}{original references}
\Crefname{originalreference}{Original reference}{Original references}
\makeatletter
\newlabel{res:rt-lb}{{2}{source}{}{sourcepointer-2}{}}
\newlabel{res:rt-lb@cref}{{[originalreference][0][]2}{[1][0][]source}{}{}{}}
\makeatother
\title{Space lower bounds for adversarially robust missing-item algorithms using a random tape}
\author{Amit Chakrabarti, Manuel Stoeckl}
\date{doi:10.4230/LIPIcs.CCC.2024.28; CC BY-SA4.0}
\begin{document}
\maketitle
\noindent\textit{Editorial note.} Original source: Amit Chakrabarti and Manuel Stoeckl, arXiv:2310.03634v2 (2024-07-03), https://arxiv.org/abs/2310.03634v2. License: Creative Commons Attribution-ShareAlike4.0, https://creativecommons.org/licenses/by-sa/4.0/.\par\medskip
\noindent\textit{Editorial note.} Changes: selected complete author theorem/proof/context excerpts and original diagram/macros/bibliography in a standalone layout; omitted unrelated sections and author frontmatter. No mathematical corrections or additional proofs. This adapted presentation retains CC BY-SA4.0.\par\medskip
\noindent\textit{Allowed standard prerequisites (editorial).}\begin{itemize}
\item One-way avoid communication lower bound. Author Theorem lem:avoid-lb explicitly cites ChakrabartiGS22; complete quoted statement and problem definition retained.
\item Multiplicative Chernoff bounds under negative association. Author Lemma lem:chernoff-neg-assoc explicitly cites JoagDevP83; statement retained.
\end{itemize}\medskip
\ifdefined\hypertarget\hypertarget{sourcepointer-1}{}\fi
\noindent\textit{Source reference.} Original section 2.2: number-theory/sources/ccc-2024-28/native/main-content.tex, source lines 615--615. Introductory overview pointer only; complete adversary and all proof details are included.\par\medskip
\ifdefined\hypertarget\hypertarget{sourcepointer-2}{}\fi
\noindent\textit{Source reference.} Original reference 2: number-theory/sources/ccc-2024-28/native/main-content.tex, source lines 267--267. Introductory weakened summary is a consequence of the delivered exact theorem; it is not used as a premise.\par\medskip
In this paper, we prove that for certain parameter regimes, \miflong
(henceforth, $\mif$) is such a problem. In the problem $\mif(n,\rl)$, the
input is a stream $\langle e_1,\ldots,e_\rl \rangle$ of $\rl$ integers, not
necessarily distinct, with each $e_i \in \{1,\ldots,n\}$, where $1 \le \rl \le
n$. The goal is as follows: having received the $i$th integer, output a number
$v$ in $\{1,\ldots,n\} \setminus \{e_1,\ldots,e_i\}$.  We will be mostly interested in the setting $\rl = o(n)$, so
the ``trivial'' upper bound on the space complexity of $\mif(n,\rl)$ is
$O(\rl \log n)$, achieved by the deterministic algorithm that simply stores
the input stream as is.

\subsection{Groundwork for Our Results}\label{subsec:models}

To state our results about $\mif$, we need to introduce some key terminology.
Notice that \mif is a {\em tracking problem}: an output is required after
reading each input.%
%
\footnote{We do not consider algorithms with a ``one-shot'' guarantee, to only
be correct at the end of the stream, because a) the adversarial setting
requires tracking output b) for \mif and most other problems the difference
in space complexity is generally small.} 
%
Thus, we view streaming algorithms as generalizations of finite state
(Moore-type) machines. An algorithm $\cA$ has a finite set of states $\Sigma$
(leading to a space cost of $\log_2|\Sigma|$), a finite input set $\cI$, and a
finite output set $\cO$.  It has a transition function $\trans \colon \Sigma
\times \cI \times \cR \to \Sigma$ indicating the state to switch to after
receiving an input, plus an output function $\gamma \colon \Sigma \times \cR
\to \cO$ indicating the output produced upon reaching a state. How the final
parameter (in $\cR$) of $\trans$ and $\gamma$ is used depends on the type of
randomness. We consider four cases, leading to four different models of
streaming computation.

\iflipics\begin{itemize}\else
\begin{itemize}[left=2ex]\fi
  %
  \item \emph{Deterministic.~} The initial state of the algorithm is a fixed
  element of $\Sigma$, and $\trans$ and $\gamma$ are deterministic (they do not
  depend on the parameter in $\cR$). 
  %
  \item \emph{Random seed.~} The initial state is drawn from a distribution
  $\cD$ over $\Sigma$, and $\trans$ and $\gamma$ are deterministic. This models
  the situation that all random bits used count towards the algorithm's space
  cost.
  %
  \item \emph{Random tape.~} The initial state is drawn from a distribution
  $\cD$ over $\Sigma$.\footnote{Requiring that this model use a fixed initial
  state could make some algorithms use one additional ``\textsc{init}''
  state.} The space $\cR$ is a sample space; when the algorithm receives an
  input $e \in \cI$ and is at state $\sigma \in \Sigma$, it chooses a random
  $\rho \in \cR$ independent of all previous choices and moves to state
  $\trans(e, \sigma, \rho)$. However, $\gamma$ is deterministic.%
  %
  \footnote{Alternatively, we could associate a \emph{distribution} of outputs
  to each state, or a function mapping (input, state) pairs to outputs. As
  these formulations are slightly more complicated to prove things with, and
  only affect the space usage of \miflong algorithms by an additive $O(\log n
  + \log \frac{1}{\delta})$ amount, we stick with the one state = one output
  convention.} %
  %
  This models the situation that the algorithm can make random decisions at
  any time, but it cannot remember past random decisions without recording
  them (which would add to its space cost). 
  %
  \item \emph{Random oracle.~} The initial state is fixed; $\cR$ is a
  sample space. A specific $R \in \cR$ is drawn at the start of the algorithm
  and stays the same over its lifetime. When the algorithm is at state
  $\sigma$ and receives input $e$, its next state is $\trans(e, \sigma, R)$.
  The output given at state $\sigma$ is $\gamma(\sigma, R)$.  This models the
  situation that random bits are essentially ``free'' to the algorithm; it can
  read from a long random string which doesn't count toward its space cost and
  which remains consistent over its lifetime. A random oracle algorithm can be
  interpreted as choosing a random deterministic algorithm, indexed by $R$,
  from some family.  
  %
\end{itemize}

These models form a rough hierarchy; they have been presented in (almost)
increasing order of power. Every $z$-bit ($2^z$-state) deterministic algorithm
can be implemented in any of the random models using $z$ bits of space; the
same holds for any $z$-bit random seed algorithm.  Every $z$-bit random tape
algorithm has a corresponding $(z + \log\rl)$-bit random oracle
algorithm---the added space cost is because for a random oracle algorithm to
emulate a random tape algorithm, it must have a way to get ``fresh''
randomness on each turn.%
%
\footnote{An alternative, which lets one express $z$-bit random tape
algorithms using a $z$-bit random oracle variant, is to assume the random
oracle algorithm has access to a clock or knows the position in the stream for free;
both are reasonable assumptions in practice.}

Streaming algorithms are also classified by the kind of correctness guarantee
they provide. Recall that we focus on ``tracking''
algorithms~\cite{BenEliezerJWY20}; they present an output after reading each
input item and this {\em entire sequence} of outputs must be correct.  Here
are three possible meanings of the statement ``algorithm $\cA$ is
$\delta$-error'' (we assume that $\cA$ handles streams of length $\rl$ with
elements in $\cI$ and has outputs in $\cO$):

\iflipics\begin{itemize}\else
\begin{itemize}[left=2ex]\fi
  %
  \item \emph{Static setting.~} For all inputs $\stream \in \cI^\rl$,
  running $\cA$ on $\stream$ produces incorrect output with probability $\le
  \delta$.
  %
  \item \emph{Adversarial setting.~} For all (computationally unbounded)
  adaptive adversaries $\alpha$ (i.e., for all functions
  $\alpha \colon \cO^\star \to \cI$),\footnote{By the minimax
  theorem, it suffices to consider deterministic adversaries.} running $\cA$
  against $\alpha$ will produce incorrect output with probability $\le
  \delta$.
  %
  \item \emph{Pseudo-deterministic setting.~} There exists a canonical output
  function $f \colon \cI^\star \rightarrow \cO$ producing all correct outputs
  so that, for each $\stream \in \cI^\rl$, $\cA(\stream)$ fails to output
  $f(\stream)$ with probability $\le \delta$.
  %
\end{itemize}

Algorithms for the static setting are called ``classic'' streaming algorithms;
ones for the adversarial setting are called ``adversarially robust'' streaming
algorithms. All pseudo-deterministic algorithms are adversarially robust, and
all adversarially robust algorithms are also classic.

As a consequence of Newman's theorem~\cite{Newman91}, any random oracle or
random tape algorithm in the static setting with error $\delta$ can be
emulated using a random seed algorithm with only $\epsilon$ increase in error
and an additional $O(\log \rl + \log \log |\cI| + \log \frac{1}{\epsilon
\delta})$ bits of space. However, the resulting algorithm is non-constructive.


\section{Preliminaries}\label{sec:prelim}

\mypara{Notation} Throughout this paper, $\log x = \log_2 x$, while $\ln x =
\log_e x$. The set $\NN$ consists of all positive integers; $[k] :=
\{1,2,\ldots,k\}$; and $[a,b)$ is a half open interval of real numbers. For a
condition or event $E$, the symbol $\indic_{E}$ takes the value $1$ if $E$
occurs and $0$ otherwise. The sequence (stream) obtained by concatenating
sequences $a$ and $b$, in that order, is denoted $a \concat b$. For a set $S$
of elements in a totally ordered universe, $\sort(S)$ denotes the sequence of
elements of $S$ in increasing order; $\binom{S}{k}$ is the set of $k$-element
subsets of $S$; and $\SSD{S}{k} = \{\sort(Y) : Y \in \binom{S}{k}\}$. We
sometimes extend set-theoretic notation to vectors and sequences; e.g., for $y
\in [n]^t$, write $y \subseteq S$ to mean that $\forall i \in [t]: y_i \in S$.
For a set $X$, $\triangle[X]$ denotes the set of probability distributions
over $X$, while $A \in_R X$ indicates that $A$ is chosen uniformly at random
from $X$. When naming probability distributions, we will use either
calligraphic letters (e.g., $\cA,\cD$) or the letter $\mu$; $A \sim \mu$ means
that $A$ is drawn from the distribution $\mu$.

As outlined in \Cref{subsec:rt-lb-overview}, $\avoid(m,a,b)$ is a one-way
communication problem, wherein player Alice has a set $A \in
\smash{\binom{[m]}{a}}$, and should send a short message to player Bob, who
should use the message to output a set $B \in \smash{\binom{[m]}{b}}$ that is
disjoint from $A$. In the randomized $\delta$-error setting, this disjointness
should hold with probability $\ge 1-\delta$.

\begin{theorem}[\avoid communication lower bound, from \cite{ChakrabartiGS22}]\label{lem:avoid-lb}
  Suppose there exists a randomized protocol for $\avoid(m,a,b)$, in which
  Alice communicates $\le K$ bits, that is $\delta$-error either on a
  worst-case input or when Alice's input is chosen uniformly at random from
  $\binom{[t]}{a}$. Then
  \begin{align*}
    K \ge \frac{a b}{t \ln 2} + \log(1-\delta) \,.
  \end{align*}
\end{theorem}

Using \Cref{lem:avoid-lb}, it is straightforward to derive the following lower
bound for robust algorithms for $\mif(n,\rl)$, as was done in
\cite{Stoeckl23}. The proof is short, yet instructive, so we outline it here.

\begin{theorem}[Adversarially robust random oracle lower bound, from \cite{Stoeckl23}]\label{lem:ext-robust-lb}
  If there exists a $\delta$-error adversarially robust random oracle
  algorithm for $\mif(n,\rl)$ using $z$ bits of space, then
  \begin{align*}
    z \ge \frac{\rl^2}{4 n \ln 2} + \log(1-\delta) \,.
  \end{align*}
\end{theorem}

\begin{lemma}[Multiplicative Azuma's inequality]\label{lem:azumanoff}
  Let $X_1,\ldots,X_t$ be $[0,1]$ random variables, and $\alpha \ge 0$. If, for all $i \in [t]$, $\EE[X_i \mid X_1,\ldots,X_{i-1}] \le p_i$, then
  \begin{align*}
    \Pr\left[ \sum_{i=1}^{t} X_i \ge (1 + \alpha) \sum_{i =1}^{t} p_i \right] 
    \le \exp\left( - ((1+\alpha) \ln (1+\alpha) - \alpha) \sum_{i =1}^{t} p_i\right) 
    \le \exp\left( - \frac{\alpha^2}{2 + \alpha} \sum_{i =1}^{t} p_i\right) \,.
  \end{align*}
  On the other hand, if for all $i$, $\EE[X_i \mid X_1,\ldots,X_{i-1}] \ge p_i$, then
  \begin{align*}
    \Pr\left[ \sum_{i=1}^{t} X_i \le (1 - \alpha) \sum_{i =1}^{t} p_i \right]
    \le \exp\left( - ((1-\alpha) \ln (1-\alpha) + \alpha) \sum_{i =1}^{t} p_i\right) 
    \le \exp\left( - \frac{\alpha^2}{2} \sum_{i =1}^{t} p_i\right) \,.
  \end{align*}
\end{lemma}

\noindent In contrast to the above, the ``usual'' form of Azuma's inequality uses a
martingale presentation and gives an additive-type bound.

\begin{lemma}[Chernoff bound with negative association, from \cite{JoagDevP83}]\label{lem:chernoff-neg-assoc} The standard multiplicative Chernoff bounds work with negatively associated random variables.
\end{lemma}

\begin{proof}[Proof of \Cref{lem:azumanoff}]
  The proof is modeled off that in \cite{Stoeckl23}, which only addresses
  one direction. It is a straightforward blend of standard proofs of the Chernoff bound and of Azuma's inequality.
%   uses a mixture of the standard approaches for the Chernoff bound and for Azuma's inequality.
%   We only prove this elementary lemma since it's hard to find a source for. The proof is modeled off that in \cite{Stoeckl23}, which only addresses one direction.

  First, the $\ge$ direction. Choose, with foresight, $z = \ln(1 + \alpha)$.
  \begin{align*}
    \Pr&\left[ \sum_{i=1}^{t} X_i \ge (1 + \alpha) \sum_{i =1}^{t} p_i \right] \\
      &= \Pr\left[ \exp( z \sum_{i=1}^{t} X_i) \ge \exp(z (1 + \alpha) \sum_{i =1}^{t} p_i) \right] \\
      &\le \frac{\EE \exp( z \sum_{i=1}^{t} X_i)}{\exp(z (1 + \alpha) \sum_{i =1}^{t} p_i)} \\
      &\le \frac{\EE[e^{z X_1} \EE[e^{z X_2} \ldots \EE[e^{z X_t} | X_1,\ldots,X_{t-1}] \ldots | X_1] ]}{\exp(z (1 + \alpha) \sum_{i =1}^{t} p_i)} \,.
  \end{align*}
  The innermost term $\EE[e^{z X_t} | X_1,\ldots,X_{t-1}]$ is, by convexity of $e^z$, $\le p_t e^z + (1-p_t) \le e^{p_t (e^z - 1)}$; after applying this upper bound, we can factor it out and bound the $X_{t-1}$ term, and so on. Thus we continue the chain of inequalities to get:
  \begin{align*}
    \le \frac{\exp\left((e^z - 1)\sum_{i=1}^{t} p_i\right)}{\exp\left(z (1 + \alpha) \sum_{i =1}^{t} p_i\right)} = \exp\left(- ((1+\alpha) \ln (1+\alpha) - \alpha) \sum_{i =1}^{t} p_i\right) \,.
  \end{align*}
  
  For the other direction, set $z = \ln(1-\alpha)$, which is $< 0$.
  This time, $\EE[e^{z X_t} | X_1,\ldots,X_{t-1}] \le p_t e^z + (1-p_t)$ because $p_t$ is a \emph{lower bound} for $\EE[X_t | X_1,\ldots,X_{t-1}]$, and $z$ is negative. That $p_t e^z + (1-p_t) \le e^{p_t (e^z - 1)}$ still holds for negative $z$, so:
  \begin{align*}
    \Pr\left[ \sum_{i=1}^{t} X_i \le (1 - \alpha) \sum_{i =1}^{t} p_i \right] &= \Pr\left[ \exp(z\sum_{i=1}^{t} X_i) \ge \exp(z (1 - \alpha) \sum_{i =1}^{t} p_i) \right] \\
      &\le \ldots \le \frac{\exp\left((e^z - 1)\sum_{i=1}^{t} p_i\right)}{\exp\left(z (1 - \alpha) \sum_{i =1}^{t} p_i\right)} \\
      &= \exp\left(- ((1-\alpha) \ln (1-\alpha) + \alpha) \sum_{i =1}^{t} p_i\right) \,. \qedhere
  \end{align*}
\end{proof}

\begin{proof}
  Such an algorithm $\cA$ yields a $z$-bit $\delta$-error randomized protocol
  for $\avoid(n, \ceil{\rl/2}, \floor{\rl/2}+1)$ wherein Alice feeds her set
  $A$ into $\cA$, sends the state of $\cA$ to Bob, and Bob extracts set $B$ by
  using the ``echo'' adversarial strategy, i.e., repeatedly asking $\cA$ for
  an output item and feeding that item back as the next input. The result now
  follows by appealing to \Cref{lem:avoid-lb}.
\end{proof}

We also note the following simple lower bound.

\begin{lemma}\label{lem:rt-triv-logr-lb}
  For every $\delta < 1$, if there exists a $\delta$-error random tape
  algorithm for $\mif(n,\rl)$ using $z$ bits of space, then $z \ge
  \log(\rl+1)$.
\end{lemma}
\begin{proof}
  Each state of a random tape streaming algorithm $\cA$ has a unique
  associated output value. If $z < \log(\rl + 1)$, then $\cA$ has at most
  $\rl$ states. Let $H$ be the set of outputs associated with these states; so
  $|H| \le \rl$. When sent a stream containing each element of $H$, $\cA$ will
  fail with probability $1$ because every output it could make is wrong.
\end{proof}

By the remarks in \Cref{subsec:models} following the definitions of the models of
computation, \Cref{lem:ext-robust-lb} also applies to random tape and random
seed algorithms and \Cref{lem:rt-triv-logr-lb} also applies to random seed
algorithms.

\section{The Random Tape Lower Bound}\label{sec:rt-lb}

This section presents our first and perhaps most important lower bound, of
which \Cref{res:rt-lb} is a consequence. We shall carry out the proof plan
outlined in \Cref{subsec:rt-lb-overview}, designing a recursive adversary to
foil a given random-tape \mif algorithm $\cA$ that runs in $z$ bits of space.
Correspondingly, our lower bound proof will be inductive.

\subsection{Setup and Base Case}

Recall that the adversary organizes the $\rl$-length input to be fed into
$\cA$ as a random prefix of length $\rl/2$ followed by another $\rl/2$ inputs
divided into several phases, each consisting of $t = O(\rl/z)$ inputs. The
adversary's eventual goal is to identify a particular phase and a
corresponding sub-universe $W \subseteq [n]$ so that $\cA$, when suitably
conditioned and restricted to that phase, yields a sub-algorithm $\cB$ that is
good for \mif for inputs from $W$. However, we will need to generalize the notion of a
``good'' \mif algorithm, because this sub-algorithm might produce outputs outside of $W$,
even when fed inputs from $W$. To aid our analysis, we will make $\cB$ abort
anytime it would have produced an output outside of $W$. In what follows, it
will be important to maintain a distinction between these aborts and actual
mistakes.

\begin{definition}\label{def:mif-rt-err-complexity}
  An algorithm $\cA$ for $\mif(n,\rl)$ can fail in either of two ways. It may
  make an incorrect output, or {\em mistake}, if outputs an element in $[n]$
  that {\em is} in its input stream (i.e., not missing). It may also {\em
  abort}, by outputting a special value $\bot$ (where $\bot \notin [n]$)
  and stopping its run.
 
  For integers $n,\rl,z$ with $1 \le \rl < n$, and $\gamma \in [0,1]$, let
  $\Algs(n,\rl,\gamma,z)$ be the set of all $z$-bit \emph{random tape}
  algorithms for $\mif(n,\rl)$ which on \emph{any} adversary abort
  with probability $\le \gamma$. Define
  \begin{align*}
    \Mstk(n,\rl,\gamma,z) := \min_{\cA \,\in\, \Algs(n,\rl,\gamma,z)} 
      \delta_{\max}(\cA, n,\rl) \,,
  \end{align*}
  where $\delta_{\max}(\cA, n,\rl)$ is the maximum probability, over all
  possible adversaries, that $\cA$ makes a mistake. As a consequence of the
  definition, $\Mstk(n,\rl,\gamma,z)$ is non-increasing in $\gamma$ and $z$.
\end{definition}

We shall establish \Cref{res:rt-lb} (concretely, \Cref{thm:rt-lb-intro}) using
a proof by induction. The base case is straightforward and handled by the
following lemma.

\begin{lemma}[Base case]\label{lem:mif-rt-lb-base-case}
  If a random tape adversarially robust algorithm $\mif(n,\rl)$ uses at most
  $\rl^2/(16 n \ln 2)$ bits of space and aborts with probability $\le \frac12$,
  then it makes a mistake with probability $\ge \frac14$. Equivalently,
  \begin{align}
    \Mstk(n,\rl,\gamma,z) 
    \ge \frac{1}{4} \indic_{z \le \rl^2 / (16 n \ln 2)} \indic_{\gamma \le 1/2} \,.
    \label{eq:mif-rt-lb-base-eq}
  \end{align}
\end{lemma}
%
\begin{proof}%[Proof of \Cref{lem:mif-rt-lb-base-case}]
  Suppose that $\gamma \le \frac12$. Let $\cA \in \Algs(n,\rl,\gamma,z)$ be an
  algorithm with mistake probability $\delta \le \frac14$. Blurring the
  distinction between aborts and mistakes, and applying the space lower bound
  for random {\em oracle} algorithms (\Cref{lem:ext-robust-lb}), we obtain
  \begin{align*}
    z \ge \frac{\rl^2}{4 n \ln 2} + \log(1 - \gamma - \delta) 
    \ge \frac{\rl^2}{4 n \ln 2} - 2 > \frac{\rl^2}{16 n \ln 2} \,.
  \end{align*}
  The latter inequality holds because $z \ge 1$ (trivially) and we have
  $\max(1, x-2) > x/4$. Taking the contrapositive, if $z \le \rl^2 / (16 n \ln
  2)$, then $\delta > 1/4$. This proves \cref{eq:mif-rt-lb-base-eq}.
\end{proof}

\subsection{The Induction Step} \label{subsec:rt-ind-step}

The induction step consists of a reduction, using an adaptive adversary
described in \Cref{alg:rt-adversary-step}, to prove a lower bound on the
mistake probability. This is formalized in the next lemma and the rest of
\Cref{subsec:rt-ind-step} is devoted to its proof.

\begin{lemma}[Induction lemma]\label{lem:mif-rt-lb-step}
  Let $1 \le \rl < n$ and $z$ be integers. Define, matching definitions in \Cref{alg:rt-adversary-step},
  \begin{align}
    w := 2 \floor{32 \frac{z n}{\rl}} \qquad \text{and} \qquad 
    t := \floor{\frac{\rl}{64 z}} \,.
    \label{eq:w-t-def}
  \end{align}
  %
  If $z \ge 8$ and $t < w$, then:
  % note: the extra 1/2 k error is _uniform_. Also have some bias from
  % probabilities...
  \begin{align}
    \Mstk\Big(n,\rl,\frac{1}{2},z\Big) \ge \min\left( \frac{\rl}{2^7 n k}, \,
      \frac{1}{4} \Mstk\Big(w,t,\frac{1}{2},z\Big) \right) \,. 
      \label{eq:mif-rt-lb-step-err}
  \end{align}
\end{lemma}

\mypara{The Initial Random Prefix} The adversary begins by sending $\cA$ a
uniformly random sequence $X \in_R \SSD{[n]}{q}$, i.e., a sorted sequence of
$q$ distinct elements of $[n]$; we'll eventually use $q = \ceil{\rl/2}$.  Let
$F$ be the {\em random function} where, for $x \in \SSD{[n]}{q}$, $F(x)$ is
the random state reached by $\cA$ upon processing $x$, starting at its initial
state; note that $F$ is determined by the transition function $T$ of $\cA$
(see~\Cref{subsec:models}).  Let $\Sigma$ be the set of states of $\cA$, so
$|\Sigma| = 2^z$. For each $\sigma \in \Sigma$, define
\begin{align}
  H_{\sigma} := \left\{ i\in[n]:\, \Pr[i\in X \mid F(X) = \sigma] \le \frac{q}{4n} \right\} \,,
  \label{eq:H-def}
\end{align}
which we can think of as the set of inputs that are ``unlikely'' to have been
seen given that $\cA$ has reached $\sigma$.  A key part of our proof of
\Cref{lem:mif-rt-lb-step} is the following lemma, which says that if $\Sigma$
is small, then $\cA$ doesn't have enough space to mark too many inputs as
unlikely, so $H_{F(X)}$ is likely to be small. The lemma can be seen as a
smoothed variant of the communication lower bound for $\avoid$.

\begin{lemma}\label{lem:forward-avoid}
  Let $\Sigma$ be a set with $|\Sigma| \le 2^z$, let $1 \le q \le n$ be
  integers, let $F$ be a random function that maps each sequence in
  $\SSD{[n]}{q}$ to a random element of $\Sigma$, and let $X \in_R
  \SSD{[n]}{q}$, chosen independently of $F$. For each $\sigma\in\Sigma$,
  define $H_\sigma$ as in \cref{eq:H-def}. Then, for all
  $\alpha\in\left(0,1\right)$,
  \begin{align*}
    \Pr\left[\left|H_{F(X)}\right|\ge\hat{w}\right] \le \alpha
    \qquad\text{where}\qquad
    \hat{w} := \left\lceil \frac{2\ln2}{1-\ln2} \frac{z+1+\log\frac{1}{\alpha}}{q}\, n\right\rceil \,.
  \end{align*}
\end{lemma}
%
\begin{proof}%[Proof of \Cref{lem:forward-avoid}]
   Consider a specific $\sigma\in\Sigma$. By linearity of expectation:
  \begin{align*}
    \EE\left[\sum_{i\in H_{\sigma}}\indic_{i \in X}\Bigm| F\left(X\right)=\sigma\right]=\sum_{i\in H_{\sigma}}\Pr\left[i\in X\mid F\left(X\right)=\sigma\right]\le\frac{q}{4n}\left|H_{\sigma}\right| \,.
  \end{align*}
  Then by Markov's inequality,
  \begin{align*}
    \Pr\left[\sum_{i\in H_{\sigma}}\indic_{i \in X}\ge2\cdot\frac{q}{4n}\left|H_{\sigma}\right|\Bigm| F\left(X\right)=\sigma\right]
      &\le \frac{1}{2} \\
      \text{which implies}\qquad
      \Pr\left[\sum_{i\in H_{\sigma}}\indic_{i \in X}\le\frac{q}{2n}\left|H_{\sigma}\right|\Bigm| F\left(X\right)=\sigma\right]
        &\ge \frac{1}{2} \,.
  \end{align*}
  Since $X$ is drawn uniformly at random from $\SSD{[n]}{q}$, the random
  variables $\{\indic_{i \in X}\}_{i \in [n]}$ are negatively associated, with
  $\EE \indic_{i \in X} = q/n$ for each $i \in [n]$. For any set $A
  \subseteq [n]$, we use the multiplicative Chernoff bound
  (\Cref{lem:chernoff-neg-assoc}) to bound the probability that $X$'s overlap
  with $A$ is much smaller than the expected value:
  \begin{align*}
    \Pr\left[\sum_{i\in A}\indic_{i \in X}\le\left(1-\frac{1}{2}\right)\cdot\frac{q}{n}|A|\right]
      \le\left(\frac{e^{-1/2}}{(1/2)^{1/2}}\right)^{(q/n)|A|}
      =\exp\left(-\frac{1}{2}(1-\ln 2)\frac{q}{n}|A|\right) \,.
  \end{align*}
  We now bound
  \begin{align*}
    \Pr\left[F\left(X\right)=\sigma\right]
    &\le \frac{\Pr\left[\sum_{i\in H_{\sigma}}\indic_{i \in X} \le \frac{q}{2n}|H_{\sigma}|\right]}%
      {\Pr\left[\sum_{i\in H_{\sigma}}\indic_{i \in X} \le \frac{q}{2n}|H_{\sigma}| \bigm| F(X)=\sigma\right]}
    \le 2\exp\left(-\frac{1}{2}(1-\ln 2)\frac{q}{n}|H_{\sigma}|\right) \,.
  \end{align*}
  Finally, let $B = \{\sigma\in\Sigma:\, |H_{\sigma}|\ge\hat{w}\}$. Then
  \begin{align*}
    \Pr\left[\left|H_{F(X)}\right|\ge\hat{w}\right]
    &=\sum_{\sigma\in B}\Pr\left[F\left(X\right)=\sigma\right] \\
    &\le 2^{z}\cdot2\exp\left(-\frac{1}{2}(1-\ln 2)\frac{q}{n}\hat{w}\right) \\
    &\le 2^{z}\cdot2\exp\left(-\left(z+1+\log\frac{1}{\alpha}\right)\ln2\right) 
    \le 2^{z}\cdot2\cdot2^{-\left(z+1+\log\frac{1}{\alpha}\right)}
    = \alpha \,. \qedhere
	\end{align*}
\end{proof}

\mypara{The Recursive Phases and Win-Win Argument} Let $\rho$ denote the
random state in $\Sigma$ reached by $\cA$ upon processing the random sequence
$X$. Having observed $\cA$'s outputs (transcript) in response to $X$, the
adversary can compute $\cD$, the distribution of $\rho$ conditioned on this
transcript. The adversary has $\floor{\rl/2}$ more items to send to $\cA$ and
it organizes these into phases of $t$ items each (see \cref{eq:w-t-def}).
Recall, from the discussion in \Cref{subsec:rt-lb-overview}, that the
adversary's goal is to identify a suitable sub-universe $W \subseteq [n]$ so
that, in some phase, $\cA$ can be seen as solving $\mif(|W|,t)$ in this
sub-universe.  As it chooses a suitable input sequence for each phase, the
adversary maintains the following objects to guide the choice:
\begin{itemize}
  \item the evolving transcript of $\cA$ given the inputs chosen so far;
  \item the corresponding distribution $\cD \in \triangle[\Sigma]$;
  \item a set $Q \subseteq \Sigma$ where, for each $\sigma \in Q$, the set
  $H_\sigma$ will be useful for determining $W$.
\end{itemize}
If the current set $Q$ leads to a suitable $W$, the adversary's strategy for
the next phase is to recursively run an optimal sub-adversary for
$\mif(|W|,t)$. If not, then (we shall show that) over the next phase, the
adversary will be able to significantly shrink the set $Q$.

To make these ideas more concrete, we introduce some terminology below; these
terms show up in \Cref{alg:rt-adversary-step}, which spells out the adversary's
actions precisely.
%
\begin{definition}[Divisive output sequence, Splitting adversary]\label{def:rt-defs}
  Let $\cA$, $\Sigma$, and $H_\sigma$ be as above and let $Q \subseteq
  \Sigma$.  A sequence (of ``outputs'') $y \in [n]^t$ is said to be
  \textsc{divisive} for $Q$ if $|\{\sigma \in Q : y \subseteq H_\sigma\}| \le
  \frac{1}{2} |Q|$.

  Say $\Upsilon$ is a $t$-length deterministic adversary, i.e., a function
  $\Upsilon \colon [n]^{\le t-1} \to [n]$.  For each $\sigma \in \Sigma$, let
  $\Outs(\sigma, \Upsilon)$ be the random variable in $[n]^t \cup \{\bot\}$
  that gives the output if we run $\cA$, starting at state $\sigma$, against
  the adversary $\Upsilon$,\footnote{This means that if, after processing a
  few inputs, the algorithm has output sequence $v \in [n]^\star$, its next
  input will be $\Upsilon(v)$.} with $\bot$ indicating that $\cA$ aborts. 
  %
  We say that $\Upsilon$ is $\alpha$-\textsc{splitting} for $Q$ with respect
  to a distribution $\cD \in \triangle[\Sigma]$ if
  \begin{align*}
    \Pr_{S \sim \cD}[\text{$\Outs(S, \Upsilon)$ is divisive for $Q$}] \ge \alpha \,,
  \end{align*} 
  with the convention that the value $\bot$ \emph{is} divisive.\footnote{The proof can also be made to work if one assumes $\bot$ is not divisive, but gives a weaker and more complicated result.}
\end{definition}

Notice that, knowing $\cD$ and $Q$, the adversary can determine whether a
given $\Upsilon$ is $\alpha$-splitting. This is implicit in 
\Cref{step:rt-pick-split-adv} of \Cref{alg:rt-adversary-step}.

\begin{algorithm}[!ht]
  \caption{An adversary for a random tape $\mif(n,\rl)$ algorithm.}
  \label[listing]{alg:rt-adversary-step}

  \begin{algorithmic}[1]
    \Statex \underline{\textsc{Adversary} (against algorithm $\cA$)}
    \vspace{.5\baselineskip}
    \State $w \gets 2 \floor{32 z n/\rl};~ h_{\max} \gets 32 z;~ t \gets \floor{\rl/(2 h_{\max})}$
    \State $X \gets$ a uniformly random sequence in $\SSD{[n]}{\ceil{\rl/2}}$.\label{step:rt-pick-v}
    \State \textbf{send} $X$ to $\cA$ and record the transcript of outputs \label{step:rt-send-v}
    \State compute $H_\sigma$ for each state $\sigma$ using \cref{eq:H-def}, with $q = \ceil{\rl/2}$
    \State $Q_0 \gets \{\sigma \in \Sigma : |H_\sigma| \le \frac{1}{2} w\}$\label{step:rt-q0-def}
    \For{$h$ in $1,\ldots,h_{\max}$}
      \State $\cD \gets$ distribution over $\Sigma$ conditioned on the cumulative transcript so far
      \If{$\exists$ a $t$-length deterministic adversary $\Upsilon$ that is $\frac12$-splitting for $Q_{h-1}$ w.r.t. $\cD$}\label{step:rt-pick-split-adv}
          \State \textbf{run} $\Upsilon$ against $\cA$ and gather the transcript of outputs $y \in [n]^t$ \label{step:rt-send-splitting}
          \State $Q_h \gets \{\sigma \in Q_{h-1}:\, y \subseteq H_\sigma\}$\label{step:rt-qh-def} \Comment{have a $\ge \frac12$ chance that $|Q_h| \le \frac{1}{2} |Q_{h-1}|$}
          \If{$Q_h = \emptyset$} \textbf{fail} \label{step:rt-abort-qempty} \EndIf
      \Else
        \State $W \gets \{i \in [n]: |\{\sigma \in Q_{h-1} : i \in H_\sigma\}|
        \ge \frac{1}{2} |Q_{h-1}|\}$ \label{step:rt-wkp1-selection} \Comment{will show that $|W| \le w$}
        \State $W' \gets W$ plus $w - |W|$ padding elements
        \State define algorithm $\cB$ to behave like $\cA$ conditioned on the transcript of inputs and outputs so far
        \State modify $\cB$, changing every output outside $W'$ to $\bot$ \label{step:rt-modify-abort} \Comment{thus $\cB$ will abort in these cases}
        % wlog this can be deterministic...
        \State let $\Xi$ be an adversary using inputs from $W'$, maximizing the probability that $\cB$ makes a mistake \label{step:rt-pick-subadv}
        \Statex \Comment{can be computed using brute-force search}
%         ensuring $\cB$ makes $\ge \Mstk(t,w,g,z)$ error. (Adversary can compute this by brute-force search.)
        \State \textbf{run} adversary $\Xi$, sending $t$ inputs in $W'$\label{step:rt-run-subadv}
        \State \textbf{return} \Comment{succeeded in causing $\cA$ to have a high enough error probability}
      \EndIf
    \EndFor
    \State \textbf{fail}\label{step:rt-abort-hmax}
  \end{algorithmic}
\end{algorithm}

We proceed to prove the induction lemma.

\begin{proof}[Proof of \Cref{lem:mif-rt-lb-step}]
  To prove the lower bound in \cref{eq:mif-rt-lb-step-err}, we show that when
  the adversary in \Cref{alg:rt-adversary-step} is run against a $z$-bit
  random-tape algorithm $\cA$ for $\mif(n,\rl)$ which has $\le \frac{1}{2}$
  worst-case probability of aborting, the probability that $\cA$ makes a
  mistake is at least the right hand side of \cref{eq:mif-rt-lb-step-err}.
  Note that the adversary feeds at most $\ceil{\rl/2} + t h_{\max} =
  \ceil{\rl/2} + \floor{\rl/ (2 h_{\max})} h_{\max} \le \ceil{\rl/2} +
  \floor{\rl/2} = \rl$ inputs to $\cA$.
  
  Consider a run of the adversary against $\cA$. This is a random process,
  with some of the randomness coming from the adversary's choices
  (\Cref{step:rt-pick-v}) and some coming from $\cA$'s internal randomness.
  Let $\rho$ be the state of $\cA$ after $X$ is sent. We now define a number
  of events, as follows.
  %
  \begin{itemize}
    % 'base'
    \item $B_{\unsafesc}$ occurs if $\cA$ produces an output in $[n] \setminus H_\rho$.
    \item $B_{\bigsc}$ occurs if the state $\rho$ has $|H_\rho| > \frac{1}{2} w$.
    \item $B_{\emptysc}$ occurs if the adversary fails at \Cref{step:rt-abort-qempty}.
    \item $B_{\timeoutsc}$ occurs if the adversary fails at \Cref{step:rt-abort-hmax}.
    \item $B_{\abortsc}$ occurs if $\cA$ aborts \emph{before} the adversary reaches \Cref{step:rt-run-subadv}.
%     \item $R_{\abortsc}(\Upsilon)$ occurs if $\cA$ aborts \emph{while} the adversary is executing \Cref{step:rt-run-subadv}, when $\Xi$ is picked to be $\Upsilon$.
    \item $R_{\errorsc}$ occurs if $\cA$ makes a mistake \emph{while} the adversary is executing \Cref{step:rt-run-subadv}.
  \end{itemize}
  
  We will consider each of the events listed above.
  \iflipics\begin{itemize}\else
  \begin{itemize}[wide, font=\bfseries]\fi
  %
  \item[(Event $B_{\unsafesc}$)~] 
    If $B_{\unsafesc}$ occurs, then some $i \in [n] \setminus H_\rho$ is
    output and, in view of \cref{eq:H-def}, the set $X$ from
    \Cref{step:rt-pick-v} has the property that
    \[
      \Pr[i \in X \mid \cA \text{ reaches } \rho] \ge \frac{\ceil{\rl/2}}{4n}
      \ge \frac{\rl}{8n} \,.
    \]
    Consequently, the probability that $\cA$ makes a mistake by producing an
    output from $X$ (which would be a non-missing item) is $\ge
    \Pr[B_{\unsafesc}] \cdot \rl/(8 n)$. 

    If $\Pr[B_{\unsafesc}] > 1/(16)$, we then have
    $\Mstk(n,\rl,\frac{1}{2},z) \ge \rl/(2^7 n)$, which implies
    \cref{eq:mif-rt-lb-step-err} leaving nothing more to prove. Therefore, for
    the rest of this proof we will consider the case in which
    \begin{align}
      \Pr[B_{\unsafesc}] \le \frac{1}{16} \,.
      \label{eq:b-unsafe-bound}
    \end{align}

  \item[(Event $B_{\bigsc}$)~]
    We apply \Cref{lem:forward-avoid} with $q = \ceil{\rl/2}$, $\alpha =
    1/16$, $F$ being the random function that maps each $x \in
    \SSD{[n]}{q}$ to the random state reached by $\cA$ upon processing $x$,
    starting at its initial state, and
    \begin{align*}
      \hat{w} = \ceil{\frac{2\ln2}{1-\ln2} \frac{z+1+\log 16}{\ceil{\rl/2}}\, n} 
        &\le 1 + \floor{\frac{2\ln2}{1-\ln2} \frac{z+1+\log 16}{\ceil{\rl/2}}\, n} \\
        &\le 1 + \floor{\frac{8 \ln 2}{1 - \ln 2} \frac{z + 1 + \log 16}{\rl}\, n} \\
        &\le 1 + \floor{19 \frac{z + 5}{\rl} n} 
          \le 1 + \floor{32 \frac{z n}{\rl}} = \frac{1}{2} w + 1 \,,
    \end{align*}
    since $z \ge 8$. As $\frac{1}{2} w$ is an integer,
    \begin{align}
      \Pr\left[B_{\bigsc}\right]
      = \Pr\left[|H_\rho| > \frac{1}{2} w\right] 
      = \Pr\left[|H_\rho| \ge \frac{1}{2} w + 1\right] 
      \le \Pr\left[|H_\rho| \ge \hat{w}\right] 
      \le \frac{1}{16} \,.
      \label{eq:b-big-bound}
    \end{align}
    
  \item[(Event $B_{\abort}$)~] 
    There are exactly exactly three spots in \Cref{alg:rt-adversary-step}
    where the algorithm can abort:
    \Cref{step:rt-send-v,step:rt-send-splitting,step:rt-run-subadv}, when the
    adversary is feeding it inputs.  By assumption, $\cA$'s worst case
    probability of aborting, against {\em any} adversary, is at most
    $\frac12$. Therefore, in particular,
    \begin{align}
      \Pr[B_{\abort}] \le \frac{1}{2} \,.
      \label{eq:b-abort-bound}
    \end{align}
    When the algorithm does abort, we stop running the adversary, so
    $B_{\abortsc}$, $B_{\emptysc}$, and $B_{\timeoutsc}$ are mutually exclusive.
    
  \item[(Event $B_{\emptysc}$)~]
    For this to happen, at some point we must have $Q_h = \emptyset$. At that
    point, the state $\rho$ must not be in $Q_h$, in which case either $\rho
    \notin Q_0$ or $\rho$ was filtered out of $Q_h$ on \Cref{step:rt-qh-def}.
    By the definition of $Q_0$, $\rho \notin Q_0$ iff $B_{\bigsc}$ holds. On
    the other hand, filtering $\rho$ out of $Q_h$ requires that the algorithm
    produce an output outside $H_{\rho}$, which can only happen if either
    $B_{\unsafesc}$ or $B_{\abortsc}$ occurs. Since $B_{\abortsc}$ is mutually
    exclusive with $B_{\emptysc}$, we conclude that
    \begin{align}
      B_{\emptysc} \Rightarrow B_{\unsafesc} \lor B_{\bigsc} \,. \label{eq:b-empty-bound}
    \end{align}
  
  \item[(Event $B_{\timeoutsc}$)~]
    We bound the probability that the adversary will fail using
    \Cref{step:rt-abort-hmax}. For this to
    happen, the adversary must have picked $h_{\max}$ splitting adversaries,
    but fewer than $z+1$ of them must have
    produced a divisive output. (If there is a divisive output in round $h$,
    then $|Q_h| \le \frac{1}{2} |Q_{h-1}|$; if not, then $|Q_h| \le
    |Q_{h-1}|$. Thus with $z+1$ divisive outputs, $|Q_{h_{\max}}| \le |Q_0| /
    2^{z+1} \le |\Sigma|/2^{z+1} \le \frac{1}{2} < 1$, in which case
    the adversary would have failed at \Cref{step:rt-abort-qempty} instead.)
  
    For each $h \in [h_{\max}]$, let $X_h$ be the $\{0,1\}$ indicator random
    variable for the event that a divisive output is found in the $h$th step.
    (If the $h$th step did not occur or no splitting adversary was found, set
    $X_h = 1$.\footnote{Note that this definition accounts for the cases where the
    algorithm aborts: if it aborts on $h$, by \Cref{def:rt-defs} this is interpreted
    as divisive, and the following steps do not occur, so $X_h=X_{h+1}=\ldots,X_{h_{\max}}=1$.
    Then the event $\{\sum X_h < z+1\}$ slightly overestimates the probability
    of $B_{\timeoutsc}$; it would be more accurate to have $\cA$ aborting produce 
    $X_h = \infty$.}) Since in the $h$th step, a splitting adversary for the
    distribution for the current state of the algorithm, conditioned on the
    transcript so far, is chosen, then $\EE[X_h \mid X_1,\ldots,X_{h-1}] \ge
    1/2$.
    %\footnote{While the adversary uses the state distribution
    %conditioned on the transcript so far, this is actually a bit more fine
    %grained than necessary. We suspect it is enough, in this part, to
    %condition on e.g. $X_1,\ldots,X_{h-1}$.} 
    Applying \Cref{lem:azumanoff} gives:
    \begin{align*}
      \Pr\Big[\sum_{h \in [h_{\max}]} X_h < z + 1\Big] &= \Pr\left[\sum_{h \in [h_{\max}]} X_h \le \left(1 - \left(1-\frac{2 z}{h_{\max}} \right) \right)\frac{h_{\max}}{2} \right]\\
        &\le \exp\left( - \frac{1}{2} \left(1-\frac{2 z}{h_{\max}} \right)^2 \frac{h_{\max}}{2}\right) \\
        &\le \exp\left( - \frac{1}{8} \frac{h_{\max}}{2}\right) &&\hspace{-2.5cm}\text{\color{black!50!white}since $h_{\max} = 32 z \ge 4 z$} \\
        &\le \frac{1}{8} \,. &&\hspace{-2.5cm}\text{\color{black!50!white}since $h_{\max} \ge 16 \ln 8$}
    \end{align*}
    Thus $\Pr[B_{\timeoutsc}] \le 1/8$.

    Combining this  with \cref{eq:b-unsafe-bound,eq:b-big-bound,eq:b-abort-bound,eq:b-empty-bound} gives:
    \begin{align}
      \Pr[B_{\emptysc} \lor B_{\timeoutsc} \lor B_{\abortsc}]
        &\le \Pr[B_{\unsafesc} \lor B_{\bigsc} \lor B_{\timeoutsc} \lor B_{\abortsc}] \nonumber\\
        &\le \Pr[B_{\unsafesc}] + \Pr[B_{\bigsc}] + \Pr[B_{\timeoutsc}] + \Pr[B_{\abortsc}] \nonumber\\
        &\le \frac{1}{16} + \frac{1}{16} + \frac{1}{8} + \frac{1}{2} = \frac{3}{4} \,.\label{eq:unsafe-big-fail-abort-bound}
    \end{align}
  
  \item[(Event $R_{\errorsc}$)~]
    Let $E$ be the event that the adversary executes
    \Cref{step:rt-run-subadv}. Notice that $R_{\errorsc}$ can only occur when
    $E$ occurs. Note also that
    \[
      E = \neg B_{\emptysc} \land \neg B_{\timeoutsc} \land \neg B_{\abort} \,,
    \]
    so
    \cref{eq:unsafe-big-fail-abort-bound} implies that 
    \begin{align}
      \Pr[E] \ge \frac14 \,.
      \label{eq:e-bound}
    \end{align}

    Suppose that $E$ does occur. We now make two claims: (a)~that $|W| \le w$,
    and (b)~that $\cB \in \Algs(w,t,\frac12,z)$. For the first claim, note
    that by \Cref{step:rt-q0-def}, the set $Q_0$ only contains states $\sigma
    \in \Sigma$ with $|H_\sigma| \le \frac{1}{2} w$. The same bound on
    $|H_\sigma|$ holds for all $\sigma \in Q_{h-1}$ because $Q_{h-1} \subseteq
    Q_0$, thanks to \Cref{step:rt-qh-def}.  By the definition of $W$
    (\Cref{step:rt-wkp1-selection}),
    \begin{align*}
      |W| 
      &= \left|\left\{i \in [n] :\, 
        \frac{|\{\sigma \in Q_{h-1} :\, i \in H_\sigma \}|}{|Q_{h-1}|} 
        \ge \frac{1}{2} \right\}\right| \\
      &\le \sum_{i \in [n]} \frac{2 |\{\sigma \in Q_{h-1} :\, i \in H_\sigma \}|}{|Q_{h-1}|} \\
      &= \frac{2}{|Q_{h-1}|} \sum_{\sigma \in Q_{h-1}} |H_\sigma|
      \le \frac{2}{|Q_{h-1}|} \cdot |Q_{h-1}| \cdot \frac{1}{2} w 
      = w \,,
    \end{align*}
    which proves claim~(a).
  
    For the second claim, consider the sub-algorithm $\cB'$ defined as $\cB$
    just before the modification at \Cref{step:rt-modify-abort}. Then, for a
    particular adversary $\Upsilon$, running $\Upsilon$ against $\cB$ causes
    an abort exactly when running $\Upsilon$ against $\cB'$ causes either an
    abort or an output outside $W'$. In the iteration of the \textbf{for} loop
    that caused $E$, since we reached the \textbf{else} branch, there was no
    deterministic splitting adversary with respect to $\cD$. Thus for
    \emph{any} deterministic adversary $\Upsilon$, by \Cref{def:rt-defs},
    \begin{align*}
      \Pr_{\hat{\sigma} \sim \cD}[\text{$\Outs(\hat{\sigma}, \Upsilon)$ is divisive for $Q_{h-1}$}] 
      \le \frac{1}{2} \,.
    \end{align*}  
    Let $y$ be a realization of the random variable $\Outs(\hat{\sigma},
    \Upsilon)$; note that $y \in [n]^{t} \cup \{\bot\}$. If a given $y$ is not
    divisive, then $y \in [n]^{t}$ and, for each $i \in y$,
    \begin{align*}
      |\{\sigma \in Q_{h-1}:\, i \in H_\sigma\}| 
      \ge |\{\sigma \in Q_{h-1}:\, y \subseteq H_\sigma\}| 
      \ge \frac{1}{2} |Q_{h-1}| \,,
    \end{align*}
    which implies that $i \in W$. Thus in fact $y \subseteq W$. It follows
    that the probability that running $\Upsilon$ against $\cB'$ causes an
    abort or an output outside $W$ (which is a subset of $W'$) is at most
    $\frac12$.
    %
    Viewing $\cB$ as a random-tape \mif algorithm handling input streams
    of length at most $t$, with items from $W'$, in the terminology of
    \Cref{def:mif-rt-err-complexity}, we have $\cB \in \Algs(w,t,\frac12,z)$.
    This proves claim~(b).
  
    Since $\Xi$ is picked to
    maximize the probability of $\cB$ making a mistake, we have
    $\Pr[R_{\errorsc} \mid E] \ge \Mstk(w,t,\frac12,z)$. Using
    \cref{eq:e-bound}, we obtain
    \begin{align*}
      \Pr[R_{\errorsc}] 
      \ge \Pr[E]\, \Mstk\left(w,t,\frac12,z\right)
      \ge \frac{1}{4} \Mstk\left(w,t,\frac12,z\right) \,.
    \end{align*}
  %
  \end{itemize}

  Combining this lower bound with the lower bound for the case where
  $\Pr[B_{\unsafesc}] > \frac{1}{16}$ (considered just before
  \cref{eq:b-unsafe-bound}), we obtain
  \begin{align*}
    \Mstk\left(n,\rl,\frac12,z\right)
    &\ge \min\left(\frac{\rl}{2^7 n},\, 
      \frac{1}{4} \Mstk\left(w,t,\frac12,z\right) \right) \,.\qedhere
  \end{align*}
\end{proof}

\subsection{Calculating the Lower Bound}

\begin{lemma}\label{lem:mif-rt-lb-induct}
  Let $1 \le \rl < n$. For any integer $k \ge 1$, say that $z$ is an integer satisfying $z \le \frac{1}{256} \rl^{1 / k}$. Then:
  \begin{align}
    \Mstk(n,\rl,0,z) > \min\Big(\frac{\rl}{2^{7} n}, \frac{1}{4^{k}} \indic_{z \le L}\Big) \qquad \text{where} \qquad L = \frac{1}{64} \left(\frac{\rl^{k + 1}}{n}\right)^{\frac{2}{k^2 + 3 k - 2}} \,. \label{eq:mif-rt-lb-multistep-lb}
  \end{align}
  Consequently, algorithms for MIF with $\le \min(\frac{\rl}{2^{7} n}, 4^{-k})$ error require $> L$ bits of space.
\end{lemma}

\begin{proof}[Proof of \Cref{lem:mif-rt-lb-induct}]
  Let $n_1 = n$ and $\rl_1 = \rl$, and for $i=2,\ldots,k$, set $n_i = 2\floor{32 \frac{z n_{i-1}}{\rl_{i-1}}}$ and $\rl_i = \floor{\frac{\rl_{i-1}}{64 z}}$. This matches the definitions used in \Cref{lem:mif-rt-lb-step}. As we have been promised that $z \le \frac{1}{256} \rl^{1 / k}$, we have in particular that:
  \begin{align*}
    (256 z)^{k} \le \rl \qquad \text{which implies} \qquad \rl_1 \ge \cdots \ge \rl_k \ge 256 z \,.
  \end{align*}
  By \Cref{lem:rt-triv-logr-lb}, we only need to consider the case $z \ge \log(\rl+1)$, as otherwise $\Mstk(n,\rl,0,z) = 1$. Thus, if $k \ge 2$, we have:
  \begin{align*}
    z \ge \log(\rl+1) \ge k \log(256 z) \ge k \log(256) = 16 \ge 8 \,.
  \end{align*}
  (If $k=1$, then \cref{eq:mif-rt-lb-multistep-lb} follows immediately from \Cref{lem:mif-rt-lb-base-case}.)

  Since $\Mstk(n,\rl,0,z) \ge \Mstk(n,\rl,1/2,z)$, it suffices to lower bound the case where algorithms are permitted up to $1/2$ abort probability. We will lower bound $\Mstk(n,\rl,1/2,z)$ by recursively applying
  \Cref{lem:mif-rt-lb-step} $k-1$ times, and then applying
  \Cref{lem:mif-rt-lb-base-case}. This yields:
  \begin{align*}
    \Mstk(n,\rl,1/2,z) &\ge \min\Big(\frac{\rl_1}{2^7 n_1}, \frac{1}{4} \Mstk(n_1,\rl_1,1/2,z)) \\
      &\ge \min\Big(\frac{\rl_1}{2^7 n_1}, \frac{1}{4} \min\Big(\frac{\rl_2}{2^7 n_2}, \ldots \frac{1}{4} \min\Big(\frac{\rl_{k-1}}{2^7 n_{k-1}},  \frac{1}{4} \indic_{z \le \rl_k^2 / (16 n_k \ln 2)} \Big)\ldots\Big)\Big) \\
      &\ge \min\Big(\frac{\rl_1}{2^7 n_1}, \frac{1}{4} \frac{\rl_2}{2^7 n_2},
      \ldots, \frac{1}{4^{k-1}} \frac{\rl_{k-1}}{2^7 n_{k-1}}, \frac{1}{4^k} \indic_{z \le \rl_k^2 / (16 n_k \ln 2)} \Big) \,.
  \end{align*}
  Only the first and last terms of the minimum are significant, because the terms for $i = 2,\ldots,{k-1}$ are all dominated by the first term:
  \begin{align*}
    \frac{\rl_i}{n_i} = \frac{\floor{\frac{\rl_{i-1}}{64 z} }}{2 \floor{32 \frac{z n_{i-1}}{\rl_{i-1}} } } \ge \frac{1}{2} \frac{\rl_{i-1}}{64 z} \frac{\rl_{i-1}}{64 z n_{i-1}} \ge \frac{\rl_{i-1}}{2 (64 z)^2} \frac{\rl_{i-1}}{n_{i-1}} \ge 4 \frac{\rl_{i-1}}{n_{i-1}} \,,
  \end{align*}
  where in the last step, we used the fact that $\rl_{i-1} \ge \rl_{k-2} \ge (256 z)^2$. Thus:
  \begin{align*}
    \Mstk(n,\rl,0,z) \ge \min\left(\frac{\rl}{2^{7} n k}, \frac{1}{4^k} \indic_{z \le \rl_k^2 / (16 n_k \ln 2)} \right) \,.
  \end{align*}
  
  We have almost proven \cref{eq:mif-rt-lb-multistep-lb}. It remains to lower bound
  $\indic_{z \le \rl_k^2 / (16 n_k \ln 2)}$ by $\indic_{z \le L}$ for some $L$. We do so by proving $z > \rl_k^2 / (16 n_k \ln 2)$ implies $z > L$. As a consequence of the definitions, we have:
  \begin{align*}
    % This is using \floor{\floor{\rl/q}/q} = \floor{\rl/q^2}
    \rl_i 
    =  \floor{\frac{\rl_{i-1}}{64 z} } 
    = \floor{\frac{\floor{\frac{\rl_{i-2}}{64 z}}}{64 z} } 
    = \cdots 
    =\floor{\frac{\rl}{(64 z)^{i-1}} } 
    \qquad \text{and} \qquad 
    n_i \le 64 \frac{z n_{i-1}}{\rl_{i-1}} \,,
  \end{align*}
  and thus:
  \begin{align*}
     z &> \frac{\rl_{k}^2}{16 \ln 2} \frac{1}{n_{k}} \ge \frac{\rl_{k}^2}{16} \frac{\prod_{k=1}^{k - 1} \rl_{k} }{(64 z)^{k - 1} n} \\
      & \ge \frac{\rl^{k + 1} }{16 \cdot (64 z)^{2 (k -1) + (k - 2) + (k - 3) + \cdots + 1 + 0}  (64 z)^{k - 1} n} \\
      & = \frac{\rl^{k + 1} }{16 \cdot (64 z)^{(k^2 + 3 k - 4) / 2} n} \,.
  \end{align*}
  Rearranging to put all the $z$ terms on the left gives:
  \begin{align*}
    64 z \cdot (64 z)^{\frac{k^2 + 3k - 4}{2}} \ge 64 \frac{\rl^{k+1}}{16 n} \,,
  \end{align*}
  which implies
  \begin{align*}
    z &> \frac{1}{64} \left(\frac{64 \rl^{k + 1}}{16 n}\right)^{\frac{2}{k^2 + 3 k - 2}} \ge \frac{1}{64} \left(\frac{\rl^{k + 1}}{n}\right)^{\frac{2}{k^2 + 3 k - 2}} \,. \qedhere
  \end{align*}
\end{proof}

Now, say a random tape algorithm has error probability $\le \frac{\rl}{2^{7} n}$ against any adaptive adversary, and uses $z$ bits of space. If $\rl \ge n^{2/3}$, it is easy to show that the optimal value of $k$ with which to apply \Cref{lem:mif-rt-lb-induct} is $1$. Otherwise, for all $k$, we have two cases: if $z \le \frac{1}{256} \rl^{1/k}$, then we can apply \Cref{lem:mif-rt-lb-induct}. By \Cref{lem:rt-triv-logr-lb}, $z \log(\rl + 1) \ge 1$, which implies $(256)^k \le \rl$ and thus $k \le \floor{\frac{1}{8} \log \rl}$. As the algorithm has
\begin{align*}
  \text{error} \le \frac{\rl}{2^{7} n} \le \min(\frac{\rl}{2^{7} n}, \frac{1}{2^{7} \rl^{1/2}} ) \le \min(\frac{\rl}{2^{7} n}, \frac{1}{4^{k}}) \,,
\end{align*}
by \Cref{lem:mif-rt-lb-induct} the algorithm's space usage must satisfy:
\begin{align*}
  z \ge \frac{1}{64} \left(\frac{\rl^{k + 1}}{n}\right)^{\frac{2}{k^2 + 3 k - 2}} \,.
\end{align*}
Since we either have this lower bound on $z$, or $z > \frac{1}{256} \rl^{1/k}$, it follows that for any integer $k \ge 1$:
\begin{align*}
  z \ge \frac{1}{256} \min\left(\rl^{1/k}, \left(\frac{\rl^{k + 1}}{n}\right)^{\frac{2}{k^2 + 3 k - 2}}\right)  \,.
\end{align*}
We can take the maximum over all values of $k$ to obtain:
\begin{align*}
  z \ge \frac{1}{256} \max_{k \in \NN} \min\Big(\rl^{1/k}, \left(\frac{\rl^{k + 1}}{n}\right)^{\frac{2}{k^2 + 3 k - 2}}\Big)  \,.
\end{align*}
The right hand side can be simplified with the following lemma, whose proof is mostly calculation and is deferred to \Cref{subsec:rt-lb-mech}:

\begin{lemma}\label{lem:rt-lb-calc-multiple}
  \begin{align}
    \max_{k \in \NN} \min\left(\rl^{1/k},\left(\frac{\rl^{k + 1}}{n}\right)^{\frac{2}{k^2 + 3 k - 2}}  \right) \ge  \max_{k \in \NN} \left(\frac{\rl^{k + 1}}{n}\right)^{\frac{2}{k^2 + 3 k - 2}}  \ge \rl^{\frac{15 \log \rl}{32 \log n}} \label{eq:rt-lb-max-min-lb}
  \end{align}
\end{lemma}

Summarizing, we obtain the following theorem.


\begin{theorem} %{rtlbintro}
\label{thm:rt-lb-intro}
  Random tape $\delta$-error adversarially robust algorithms for $\mif(n,\rl)$ require
  \begin{align*}
    \Omega\left(\max_{k \in \NN} \left(\frac{\rl^{k + 1}}{n}\right)^{\frac{2}{k^2 + 3 k - 2}}\right) 
    = \Omega\left(\rl^{\frac{15}{32} \log_{n} \rl}\right)
  \end{align*}
  bits of space, for $\delta \le \rl/(2^{7} n)$. \qed
\end{theorem}

\subsection{Mechanical Proofs for \texorpdfstring{\Cref{sec:rt-lb}}{Random Tape Lower Bound}}\label{subsec:rt-lb-mech}

\begin{proof}[Proof of \Cref{lem:rt-lb-calc-multiple}]
  Let $x$ be the left hand side of Eq. \ref{eq:rt-lb-max-min-lb}.
  The left branch of the $\min(\cdot,\cdot)$ terms in Eq. \ref{eq:rt-lb-max-min-lb} is actually unnecessary. For any integer $\lambda \ge 2$, say that
  \begin{align*}
    \left(\frac{\rl^{\lambda + 1}}{n}\right)^{\frac{2}{\lambda^2 + 3 \lambda - 2}} = \max_{k \in \NN} \left(\frac{\rl^{k + 1}}{n}\right)^{\frac{2}{k^2 + 3 k - 2}} \,.
  \end{align*}
  Then in particular,
  \begin{align*}
    \left(\frac{\rl^{\lambda + 1}}{n}\right)^{\frac{2}{\lambda^2 + 3 \lambda - 2}} \ge \left(\frac{\rl^{(\lambda -1) + 1}}{n}\right)^{\frac{2}{(\lambda-1)^2 + 3 (\lambda -1)- 2}} \ge     \left(\frac{\rl^{\lambda}}{n}\right)^{\frac{2}{\lambda^2 + \lambda - 4}} \,,
  \end{align*}
  which implies
  \begin{align*}
    n^{2 \lambda + 2} = n^{(\lambda^2 + 3 \lambda - 2) - (\lambda^2 + \lambda - 4)} \ge \rl^{\lambda \cdot (\lambda^2 + 3 \lambda - 2) - (\lambda + 1) \cdot (\lambda^2 + \lambda - 4)} = \rl^{\lambda^2 + \lambda + 4} \,,
  \end{align*}
  hence we have $\rl \le n^{\frac{2 \lambda +2}{\lambda^2 + \lambda + 4}}$. 

  On the other hand, we have
  \begin{align*}
    \rl^{1/\lambda} \ge \left(\frac{\rl^{\lambda + 1}}{n}\right)^{\frac{2}{\lambda^2 + 3 \lambda - 2}} \qquad \iff \qquad n^\lambda \le \rl^{(\lambda + 1) \lambda - \frac{\lambda^2+3\lambda - 2}{2}} = \rl^{\frac{\lambda^2 - \lambda + 2}{2}} \,,
  \end{align*}
  so the left branch of the $\min(\cdot,\cdot)$ in Eq. \ref{eq:rt-lb-max-min-lb} is only smaller when $\rl \ge n^{\frac{2 \lambda}{\lambda^2 - \lambda + 2}}$. As
  \begin{align*}
    \frac{2 \lambda +2}{\lambda^2 + \lambda + 4} \le \frac{2 \lambda}{\lambda^2 - \lambda + 2} \,,
  \end{align*}
  for all $\lambda \ge 1$, it follows that the left branch of the $\min(\cdot,\cdot)$ in Eq. \ref{eq:rt-lb-max-min-lb} is only smaller than the right when the entire term is not the maximum. Thus
  \begin{align*}
    x &\ge \max_{k \in \NN} \left(\frac{\rl^{k + 1}}{n}\right)^{\frac{2}{k^2 + 3 k - 2}} \,.
  \end{align*}

  To get a looser but more easily comprehensible lower bound, we note that $\max_{k \in \NN} \log \left(\frac{\rl^{k + 1}}{n}\right)^{\frac{2}{k^2 + 3 k - 2}}$ is piecewise linear and convex in $\log \rl$. Consequently,
  we can lower bound it using the convex function $C \frac{(\log \rl)^2}{\log n}$, where $C$ is the maximum value which satisfies the inequality at all ``corner points'' of $\max_{k \in \NN} \log \left(\frac{\rl^{k + 1}}{n}\right)^{\frac{2}{k^2 + 3 k - 2}}$. These corner points occur precisely at values of $\log \rl$ where, for some $k \ge 2$, we have:
  \begin{align*}
    \left(\frac{\rl^{k + 1}}{n}\right)^{\frac{2}{k^2 + 3 k - 2}} = \left(\frac{\rl^{(k-1) + 1}}{n}\right)^{\frac{2}{(k-1)^2 + 3 (k-1) - 2}} \,.
  \end{align*}
%   which is equivalent to: 
%   \begin{align*}
%     \left(\frac{k + 1}{k^2 + 3 k - 2} - \frac{k}{(k-1)^2 + 3 (k-1) - 2}\right) \log \rl = \left(\frac{1}{k^2 + 3 k - 2} -  \frac{1}{(k-1)^2 + 3 (k-1) - 2}\right) \log n
%   \end{align*}
  Rearranging this gives:
  \begin{align*}
    \frac{\log n}{\log \rl} = \frac{(k+1) ((k-1)^2 + 3 (k-1) - 2) - (k-1+1) (k^2 + 3 k - 2)}{((k-1)^2 + 3 (k-1) - 2) - (k^2 + 3 k - 2)} = \frac{k^2 + k + 4}{2k + 2} \,.
  \end{align*}
  so the corners occur at $\rl = n^{\frac{2k + 2}{k^2 + k + 4}}$; and at such $\rl$, we have
  \begin{align*}
    \log \left(\frac{\rl^{k + 1}}{n}\right)^{\frac{2}{k^2 + 3 k - 2}} 
    &= \left(\frac{2}{k^2 + 3 k - 2} \cdot ((k+1) \frac{2 k + 2}{k^2 + k + 4} - 1)\right) \log n = \frac{2}{k^2 + k + 4} \log n \\
      &= \frac{(\log \rl)^2}{\log n} \frac{2}{k^2 + k + 4} \left(\frac{\log n}{\log \rl}\right)^2 = \frac{(\log \rl)^2}{\log n} \frac{2}{k^2 + k + 4} \left(\frac{k^2 + k + 4}{2k + 2}\right)^2 \\
      %&=  \frac{(\log \rl)^2}{\log n} \frac{2 (k^2 + k + 4)}{(2k + 2)^2} 
      &= \frac{1}{2} \frac{(\log \rl)^2}{\log n} \frac{k^2 + k + 4}{(k+1)^2} \ge \frac{15}{32} \frac{(\log \rl)^2}{\log n} \,.
  \end{align*}
  The function $\frac{k^2 + k + 4}{(k+1)^2}$ has derivative $\frac{k-7}{(k+1)^3}$ and is minimized exactly at $k = 7$, where it has value $\frac{15}{16}$. Consequently, the value $C = \frac{15}{32}$ is the best possible.
\end{proof}

\newcommand{\etalchar}[1]{$^{#1}$}
\begin{thebibliography}{GGMW20}

\bibitem[ABJ{\etalchar{+}}22]{AjtaiBJSSWZ22}
Mikl\'os Ajtai, Vladimir Braverman, T.S. Jayram, Sandeep Silwal, Alec Sun,
  David~P. Woodruff, and Samson Zhou.
\newblock The white-box adversarial data stream model.
\newblock In {\em Proc. 41st ACM Symposium on Principles of Database Systems},
  pages 15--27, 2022.

\bibitem[ACGS23]{AssadiCGS23}
Sepehr Assadi, Amit Chakrabarti, Prantar Ghosh, and Manuel Stoeckl.
\newblock Coloring in graph streams via deterministic and adversarially robust
  algorithms.
\newblock In {\em Proc. 42nd ACM Symposium on Principles of Database Systems},
  pages 141--153, 2023.

\bibitem[BEEO22]{BenEliezerEO22}
Omri Ben-Eliezer, Talya Eden, and Krzysztof Onak.
\newblock Adversarially robust streaming via dense-sparse trade-offs.
\newblock In {\em Symposium on Simplicity in Algorithms (SOSA)}, pages
  214--227, 2022.

\bibitem[BJWY20]{BenEliezerJWY20}
Omri Ben{-}Eliezer, Rajesh Jayaram, David~P. Woodruff, and Eylon Yogev.
\newblock A framework for adversarially robust streaming algorithms.
\newblock In {\em Proc. 39th ACM Symposium on Principles of Database Systems},
  page 63–80, 2020.

\bibitem[BKKS23]{BravermanKKS23}
Vladimir Braverman, Robert Krauthgamer, Aditya Krishnan, and Shay Sapir.
\newblock Lower bounds for pseudo-deterministic counting in a stream.
\newblock {\em arXiv preprint arXiv:2303.16287}, 2023.

\bibitem[BR93]{BellareR93}
Mihir Bellare and Phillip Rogaway.
\newblock Random oracles are practical: A paradigm for designing efficient
  protocols.
\newblock In {\em Proc. 1st ACM Conference on Computer and Communications
  Security}, pages 62--73, 1993.

\bibitem[BY20]{BenEliezerY20}
Omri Ben{-}Eliezer and Eylon Yogev.
\newblock The adversarial robustness of sampling.
\newblock In {\em Proc. 39th ACM Symposium on Principles of Database Systems},
  pages 49--62. {ACM}, 2020.

\bibitem[CGS22]{ChakrabartiGS22}
Amit Chakrabarti, Prantar Ghosh, and Manuel Stoeckl.
\newblock Adversarially robust coloring for graph streams.
\newblock In {\em Proc. 13th Conference on Innovations in Theoretical Computer
  Science}, pages 37:1--37:23, 2022.

\bibitem[Fei19]{Feige19}
Uriel Feige.
\newblock A randomized strategy in the mirror game.
\newblock {\em arXiv preprint arXiv:1901.07809}, 2019.

\bibitem[GGMW20]{GoldwasserGMW20}
Shafi Goldwasser, Ofer Grossman, Sidhanth Mohanty, and David~P. Woodruff.
\newblock {Pseudo-Deterministic Streaming}.
\newblock In {\em Proc. 20th Conference on Innovations in Theoretical Computer
  Science}, volume 151, pages 79:1--79:25, 2020.

\bibitem[GGS23]{GrossmanGS23}
Ofer Grossman, Meghal Gupta, and Mark Sellke.
\newblock Tight space lower bound for pseudo-deterministic approximate
  counting.
\newblock {\em arXiv preprint arXiv:2304.01438}, 2023.

\bibitem[GS18]{GargS18}
Sumegha Garg and Jon Schneider.
\newblock {The Space Complexity of Mirror Games}.
\newblock In {\em Proc. 10th Conference on Innovations in Theoretical Computer
  Science}, pages 36:1--36:14, 2018.

\bibitem[HKM{\etalchar{+}}20]{HassidimKMMS20}
Avinatan Hassidim, Haim Kaplan, Yishay Mansour, Yossi Matias, and Uri Stemmer.
\newblock Adversarially robust streaming algorithms via differential privacy.
\newblock In {\em Advances in Neural Information Processing Systems 33: Annual
  Conference on Neural Information Processing Systems 2020, NeurIPS 2020,
  December 6-12, 2020, virtual}, 2020.

\bibitem[HW13]{HardtW13}
Moritz Hardt and David~P. Woodruff.
\newblock How robust are linear sketches to adaptive inputs?
\newblock In {\em Proc. 45th Annual ACM Symposium on the Theory of Computing},
  pages 121--130, 2013.

\bibitem[Ind06]{Indyk06}
Piotr Indyk.
\newblock Stable distributions, pseudorandom generators, embeddings, and data
  stream computation.
\newblock {\em J. ACM}, 53(3):307--323, 2006.

\bibitem[JP83]{JoagDevP83}
Kumar {Joag-Dev} and Frank Proschan.
\newblock Negative association of random variables, with applications.
\newblock {\em Ann. Stat.}, 11(1):286--295, 1983.

\bibitem[JW23]{JayaramW23}
Rajesh Jayaram and David~P Woodruff.
\newblock Towards optimal moment estimation in streaming and distributed
  models.
\newblock {\em ACM Trans. Alg.}, 19(3):1--35, 2023.

\bibitem[KMNS21]{KaplanMNS21}
Haim Kaplan, Yishay Mansour, Kobbi Nissim, and Uri Stemmer.
\newblock Separating adaptive streaming from oblivious streaming using the
  bounded storage model.
\newblock In {\em Advances in Cryptology - {CRYPTO} 2021 - 41st Annual
  International Cryptology Conference, {CRYPTO} 2021, Virtual Event, August
  16-20, 2021, Proceedings, Part {III}}, volume 12827 of {\em Lecture Notes in
  Computer Science}, pages 94--121. Springer, 2021.

\bibitem[MN22a]{MagenN22}
Roey Magen and Moni Naor.
\newblock Mirror games against an open book player.
\newblock In {\em 11th International Conference on Fun with Algorithms (FUN
  2022)}, volume 226, pages 20:1--20:12, 2022.

\bibitem[MN22b]{MenuhinN22}
Boaz Menuhin and Moni Naor.
\newblock Keep that card in mind: Card guessing with limited memory.
\newblock In {\em Proc. 13th Conference on Innovations in Theoretical Computer
  Science}, pages 107:1--107:28, 2022.

\bibitem[New91]{Newman91}
Ilan Newman.
\newblock Private vs.~common random bits in communication complexity.
\newblock {\em Inform. Process. Lett.}, 39(2):67--71, 1991.

\bibitem[Nis90]{Nisan90}
Noam Nisan.
\newblock Pseudorandom generators for space-bounded computation.
\newblock In {\em Proc. 22nd Annual ACM Symposium on the Theory of Computing},
  pages 204--212, 1990.

\bibitem[Nis93]{Nisan93B}
Noam Nisan.
\newblock On read once vs. multiple access to randomness in logspace.
\newblock {\em Theoretical Computer Science}, 107(1):135--144, 1993.

\bibitem[Sto23]{Stoeckl23}
Manuel Stoeckl.
\newblock Streaming algorithms for the missing item finding problem.
\newblock In {\em Proc. 34th Annual ACM-SIAM Symposium on Discrete Algorithms},
  pages 793--818, 2023.
\newblock Full version at arXiv:2211.05170v1.

\bibitem[WZ22]{WoodruffZ22}
David~P. Woodruff and Samson Zhou.
\newblock Tight bounds for adversarially robust streams and sliding windows via
  difference estimators.
\newblock In {\em Proc. 62nd Annual IEEE Symposium on Foundations of Computer
  Science}, pages 1183--1196, 2022.

\end{thebibliography}
\end{document}
